% Shared two-arm terminal proof at roots of minimum one and two.
\subsection{A two-arm terminal at roots of minimum one and two}
\label{sec:odd-two:roots}

Use the original model of Section~\ref{sec:odd-common-calculus} and put
\(h=17/32\). All lengths below are positive integers, all ordered lists
are finite, and all operators use the same original probability measure.

\begin{lemma}[Transporting a complete constant]
\label{lem:odd-two:constant}
Suppose an actual positive message satisfies
\[
 \log H\le a+B(V+c\one)-\tfrac3{10}(u+\xi Bu)^2,
 \qquad \xi\in\{0,1\},\quad c\in\mathbb R.
\]
Its two available fourth certificates are
\begin{align}
 J_f&=4(a+\E BV)+\tfrac{17}{4}\|V+c\one\|^2,
       \label{eq:odd-two:full-certificate}\\
 J_c&=4\{a+\E B(V-\xi cu)\}
       +\tfrac{17}{4}\|V-\xi cu\|^2+\tfrac{10}{3}c^2.
       \label{eq:odd-two:curved-certificate}
\end{align}
Each bounds \(\log\E H^4\). The case \(\xi=0\) uses the actual
\(p\) curvature, and \(\xi=1\) uses the actual \(W_2\) curvature.
\end{lemma}
\begin{proof}
For \(w\in[-1,1)\), \(\log(1-w)\le-w-3w^2/10\).
Its derivative proof is the unit-arm proof applied to \(w\); the
endpoint \(-1\) is paid by \(\log2<7/10\). Both \(w=u\) and
\(w=u+Bu=1-W_2\) have this range. The full certificate follows by
discarding the nonpositive square and applying
\eqref{eq:oc:signed-moments}. For the second use the exact inequality
\[
 4cu-\tfrac65(u+\xi Bu)^2
 \le\tfrac{10}{3}c^2-4\xi cBu
\]
before that same signed-mean bound. In particular no mean or
constant-vector cross term is declared zero.
\end{proof}

\begin{lemma}[The terminal packet actually needed]
\label{lem:odd-two:packet}
For \(t\ge3\), set \(m=t-2\), \(G_m=e+zf\) at
\(z=i\sqrt x\), and
\[
 \gamma=\beta_{t-1}=\int f\,d\nu,\quad
 \Gamma=\beta_t+\gamma=\int(e+f)\,d\nu,\quad
 R=(\gamma+x e+zx f)u.
\]
The actual terminal \(F_t=\bar f_2W_t\) satisfies
\(F_t\le e^{R-\Gamma}\). At \(t=3\) the same envelope also
holds for the distinct terminal \(h_2f_3\).
Put \(D_j=\E B^jR\), \(\sigma=\E R^2\),
\(\delta=D_0+\eta\sigma\), and \(\epsilon=\delta+D_1\), where
\(\eta=3/2\) at \(t\ge4\) and \(\eta=8/7\) at \(t=3\).
Then
\begin{equation}\label{eq:odd-two:packet}
\begin{gathered}
 TF_t\le e^{-\Gamma}(q+BR+\delta),\qquad
 \Gamma\ge d,\quad0\le\gamma\le d,\quad\|R\|\le\tfrac45\sqrt d,\\
 -\delta_2\le\delta\le\tfrac7{12}d,\qquad
 -\tfrac{13}{16}d\le\epsilon\le0,\qquad D_1\le0,\\
 \Gamma-2\delta_+\ge\tfrac35d,\qquad
 \Gamma-(1-d)\delta\ge\tfrac45d,\qquad
 \gamma_j:=\beta_j-D_{j-1}\ge hd\quad(j\ge3).
\end{gathered}
\end{equation}
The first-row statement also holds for the stated \(h_2f_3\) choice.
\end{lemma}
\begin{proof}
Retain \(\beta_{t-1}W_1\) in the positive \(A\)-image defining
\(W_t\). Adding \(\bar f_2\) gives
\(\bar f_2+W_t\le2-\Gamma+R\); positive AM--GM proves the
envelope. At \(t=3\) both actual sums
\(\bar f_2+W_3\) and \(h_2+f_3\) equal \(2-d-B^2u\).
The exact moments are
\begin{equation}\label{eq:odd-two:packet-moments}
\begin{aligned}
 \sigma&=\gamma^2d+2\gamma\int x e\,d\nu
                +\int x^2(e^2+xf^2)\,d\nu,\\
 D_0&=-\int xf\,d\nu,&
 D_1&=-\gamma d-\int xe\,d\nu,\\
 D_2&=\int x^2f\,d\nu,&
 D_3&=\gamma\delta_2+\int x^2e\,d\nu.
\end{aligned}
\end{equation}
For \(m\ge2\), \(1-x\le e,f\le1\). The row bound gives
\(|R|\le d+\sqrt{\delta_2+\delta_3}<1\); hence
\(e^R\le1+R+3R^2/4\) and the positive first row follows.
The positive remainder coefficient is at most \(e-2<3/4\);
the negative remainder is at most \(1/2\).
Moreover
\[
 \sigma/d\le\tfrac19+\tfrac{2\theta}{3}+\theta^2+\theta^3
 <\tfrac58<\tfrac{16}{25}.
\]

For fixed actual \(e,f\), \(\epsilon\) is convex in
\(\gamma\in[0,1/3]\). At zero, its two negative linear terms
absorb the positive squares. At \(\gamma=1/3\) it is at most
\[
 -d/6+\int[-x+(5/2)x^2+(3/2)x^3]\,d\nu\le0.
\]
Here maximize the actual \(e\)-term at one and the \(f\)-term
at \(1-x\); the latter derivative is negative. The displayed
convex cubic has endpoints zero and \(7007/65536<1/6\).
Conversely \(\sigma\ge\gamma^2d+2\gamma\int xe\,d\nu\)
gives
\[
 \epsilon/d\ge-2\theta+(3\theta-1)\gamma+(3/2)\gamma^2
 \ge-13/16.
\]
Here \(3\gamma-1\le0\) and both first integrals are at most
\(\delta_2\le\theta d\).

We prove the joint scalar bound before using it to bound \(\delta\).
When \(\delta\le0\), \(\Gamma\ge2\kappa d>5d/6\).
Otherwise consider \(\Gamma-2\delta\), retaining
\(\gamma=\int f\,d\nu\). If \(m\) is even, \(e=f\). With
\(X=\int xf\,d\nu\), \(Y=\int x^2(1+x)f^2\,d\nu\),
\[
 \Gamma-2\delta=2\gamma-3\gamma^2d+(2-6\gamma)X-3Y,
 \qquad Y\le\theta(1+\theta)X.
\]
If \(2-6\gamma-3\theta(1+\theta)\ge0\), the result is at
least \(95d/96\). Otherwise use \(X\le\theta\gamma\) and
\(a=\gamma/d\in[\kappa,1]\). The lower bound divided by \(d\) is
\[
 C a-\tfrac{55}{48}a^2,\qquad
 C=2+2\theta-3\theta^2(1+\theta)>21/10.
\]
Its two endpoint lower bounds are \(229/240\) and
\(207157/245760\), both exceeding \(5/6\).
If \(m\ge3\) is odd, \(e=1-xf\) and
\[
 \sigma=\gamma^2d+2\gamma\delta_2-2\gamma\int x^2f\,d\nu
 +\delta_3-\int x^3 f[2-(1+x)f]\,d\nu.
\]
Using \(\int xf\,d\nu\ge\kappa\delta_2\),
\(\int x^2f\,d\nu\ge\kappa\delta_3\), and
\(f[2-(1+x)f]\ge\kappa^2\) gives
\[
 \Gamma-2\delta\ge d+\gamma-3\gamma^2d+\int J_\gamma\,d\nu,
 \quad J_\gamma=(\kappa-6\gamma)x+(6\gamma\kappa-3)x^2
                         +3\kappa^2x^3.
\]
Its second derivative is at most
\(4\kappa-6+18\kappa^2\theta<0\), and its endpoint at
\(\theta\) is negative: the \(\gamma\)-coefficient is negative
and the \(\gamma=0\) value is negative. After using \(d\le1/3\), its lower bound divided
by \(d\) is
\[
 A+Ca-a^2/3,\quad
 A=1+\kappa\theta-3\theta^2+3\kappa^2\theta^3>13/16,
 \quad C=1-2\theta+2\kappa\theta^2>3/8.
\]
The two endpoint bounds exceed \(41/48\) and \(233/256\).
Thus, for every finite \(m\ge2\),
\begin{equation}\label{eq:odd-two:joint-loss}
 \Gamma-2\delta_+\ge5d/6\qquad(t\ge4).
\end{equation}
Since \(\Gamma\le2d\), this already implies the sufficient
\(\delta\le7d/12\). Also \(\delta\ge D_0\ge-\delta_2\),
and \((1-d)\delta\le2\delta_+\) proves the other scalar bound.

At \(m=1\), \((e,f,\gamma,\Gamma)=(1,0,0,d)\),
\(R=-B^2u\), and \(|R|\le\sqrt{\delta_2}<3/8\).
The finite factorial-tail bound for the remainder coefficient is
\(7281/12800<4/7\). Therefore
\[
 \delta=8\delta_3/7<d/5,\quad
 \epsilon=\delta-\delta_2,\quad D_2=0,\quad D_3=\delta_3.
\]
These values give all the stated scalar bounds.
Finally the complete four residue classes, for every \(j\ge0\), are
\[
\begin{aligned}
 D_{4j}&=-\int x^{2j+1}f\,d\nu,&
 D_{4j+1}&=-\gamma\delta_{2j+1}-\int x^{2j+1}e\,d\nu,\\
 D_{4j+2}&=\int x^{2j+2}f\,d\nu,&
 D_{4j+3}&=\gamma\delta_{2j+2}+\int x^{2j+2}e\,d\nu.
\end{aligned}
\]
The first two positive classes give
\(D_2\le\delta_3\), \(D_3\le d\delta_2+\delta_3\).
Every later positive value is at most
\((d+\theta)\delta_4<3\delta_4/4<d/16\), and each
subsequent member contracts by \(\theta^2\). Combining with the
finite \(\beta_j\) bounds proves \(\gamma_j\ge17d/32\), including
the next mean index at every connector length.
\end{proof}

\Needspace{20\baselineskip}
\begin{theorem}[The two root minima with the exact terminal \((2,t)\)]
\label{thm:odd-two:roots}
Let the root residual list be
\((r,a,b,s_1,t_1,\ldots,s_k,t_k)\), with \(r\in\{1,2\}\),
\(3\le a<b<s_1<t_1<\cdots<s_k<t_k\), and terminal list
\((2,t)\), \(t\ge3\). The directions are as follows:
\begin{longtable}{@{}ccp{.40\linewidth}p{.32\linewidth}@{}}
\caption{The two root minima and their terminal directions.}
\label{tab:odd-two:directions}\\
\toprule
\(r\)&\(L\)&Root factor&Terminal factor\\
\midrule
1&all&\(pW_a\bar f_b\prod W_{s_i}\bar f_{t_i}\)
 &\(h_2f_3\) if \(t=3\), otherwise \(\bar f_2W_t\)\\
2&1&\(f_2h_af_b\prod h_{s_i}f_{t_i}\)&\(W_2\bar f_t\)\\
2&\(\ge2\), \(a=3\)&\(W_2\bar f_3W_b\prod W_{s_i}\bar f_{t_i}\)
 &\(\bar f_2W_t\)\\
2&\(\ge2\), \(a\ge4\)&\(W_2W_a\bar f_b\prod W_{s_i}\bar f_{t_i}\)
 &\(\bar f_2W_t\)\\
\bottomrule
\end{longtable}
Every inner edge points terminal-to-root and every outer edge
root-to-receiver. Let \(H_k\) be the actual root factor times the
actual \(T^L\)-propagated terminal, with arbitrary compatible equal-pair
decorations at both roots. Put \(z_0=0\) and, for \(k\ge1\),
\(z_k=2737/3840+64(k-1)/5\). Then for every original host,
\begin{equation}\label{eq:odd-two:conclusion}
 \E H_k^4\le e^{-(6+z_k)d},\qquad
 I_q(T^MH_k)^3\le e^{-[12+2z_k-c(M)]d}\le e^{-24d/19},
\end{equation}
for every independent \(L,M\ge1\), with
\(c(1)=204/19,c(2)=10,c(3)=0,c(4)=7,c(M)=6\) for \(M\ge5\).
\end{theorem}
\begin{proof}
First omit decorations and later pairs. The actual pair envelopes are
Lemma~\ref{lem:odd-common:root-three}. Write
\(G_{b-3}=E+zC\), \(U=(xE+zxC)u\), \(\beta=\beta_{b-1}\).
The reversed pair is bounded by
\(e^{-g_r+B(-U-\tau\one)}\), with
\(g_r=3d-2\delta_2\). The exchanged pair is bounded by
\(e^{-g_e+B(U+\beta\one)}\), where
\(g_e=d+\beta_b+\beta\ge3d-2\delta_2\) and
\(\kappa d\le\beta\le d\).
At the finite boundary \(b=4\), \((E,C)=(1,0)\) and \(\beta=d\).
For a short root arm at least four, the one-pair envelope is
\(e^{-16d/5+BV_*}\), with
\(\|V_*\|\le3\sqrt d/8\), \(\E BV_*\le d/12\).
The same bounds apply in the positive direction at a unit inner edge.
They follow from the actual pair loss and the single finite mask:
\[
 \|V_*\|^2\le(101/50)\theta^3d<9d/64,
 \qquad \E BV_*\le\delta_4/(1-\theta^2)<d/12.
\]

\paragraph{The common finite root calculation.}
Use Lemma~\ref{lem:odd-two:packet}. The positive finite Duhamel majorant is
\[
 Y_L=\bar f_L+\delta\bar f_{L-1}+B^LR
                   +\sum_{j=1}^{L-1}D_j\bar f_{L-1-j}.
\]
For \(L\ge4\), the two actual final losses give
\(Y_L\le1+u+2\delta_+-2hd+B[G_nu+B^nR-hd\one]\),
where \(n=L-1\ge3\). Its next-index mean is
\(\E B(G_nu+B^nR)=-\gamma_{L+1}\le-hd\).
Before adding a root pair, the following table is exact:
\[
\begin{array}{c|c|c|c}
L&a_L&V_L&c_L\\\hline
1&-\Gamma+\delta&R&0\\
2&-\Gamma+\epsilon&u+BR&\delta\\
3&-\Gamma+\epsilon-d+D_2&(1+\delta+B)u+B^2R&\epsilon\\
\ge4&-\Gamma+2\delta_+-2hd&G_nu+B^nR&-hd
\end{array}
\]
For the actual \(p\) prefix,
\(\log(pT^LF)\le a_L+B(V_L+c_L\one)-3u^2/10\).
For the actual \(W_2\) prefix, at \(L\ge2\), replace only
\(V_L\) by \(V_L-u\), and retain \(-3(u+Bu)^2/10\).
Thus the same finite input supplies both root minima. Constants are paid
by Lemma~\ref{lem:odd-two:constant}, not discarded from this table.

The following entries record the complete payments. Here \(A\) means
\(a+\E BV\le-Ad\); \(N\) bounds the square of the final full or
transported vector divided by \(d\); \(C_0\) bounds the separate
curvature expense divided by \(d\). Each loss is \(4A-17N/4-C_0\).
\begingroup\small
\Needspace{8\baselineskip}
\begin{longtable}{@{}ccccc@{}}
\caption{Complete root certificates for the nonunit packet.
Every displayed loss exceeds six.}
\label{tab:odd-two:payments}\\
\toprule
Root and length&\(A\)&\(N\)&\(C_0\)&Loss\\\midrule
\endfirsthead
\toprule
Root and length&\(A\)&\(N\)&\(C_0\)&Loss\\\midrule
\endhead
\(pW_3\bar f_b,1\)&\(12/5\)&\(2/3\)&0&\(203/30\)\\
\(pW_3\bar f_b,2\)&\(18/5\)&\(15/8\)&0&\(1029/160\)\\
\(pW_3\bar f_b,3\)&\(179/40\)&\(15/8\)&\(15/4\)&\(989/160\)\\
\(pW_3\bar f_b,\ge4\)&\(159/40\)&\(25/16\)&\(8/3\)&\(6329/960\)\\
\(pW_a\bar f_b,1; a\ge4\)&\(247/60\)&\((47/40)^2\)&0&\(203501/19200\)\\
\(pW_a\bar f_b,2; a\ge4\)&\(76057/15360\)&\((67/40)^2\)&0&\(75673/9600\)\\
\(pW_a\bar f_b,3; a\ge4\)&\(349/60\)&4&0&\(94/15\)\\
\(pW_a\bar f_b,\ge4; a\ge4\)&\(2549/480\)&\(64/25\)&\(1/3\)&\(6017/600\)\\
\(W_2\bar f_3W_b,2\)&\(51/16\)&\(227/144\)&0&\(3485/576\)\\
\(W_2\bar f_3W_b,3\)&\(319/80\)&\(3/2\)&0&\(383/40\)\\
\(W_2\bar f_3W_b,\ge4\)&\(541/160\)&\(25/16\)&0&\(2203/320\)\\
\(W_2W_a\bar f_b,2; a\ge4\)&\(79/20\)&\((53/40)^2\)&0&\(53367/6400\)\\
\(W_2W_a\bar f_b,3; a\ge4\)&\(289/60\)&\((41/24)^2\)&0&\(79067/11520\)\\
\(W_2W_a\bar f_b,\ge4; a\ge4\)&\(2069/480\)&\(64/25\)&0&\(3817/600\)\\
\bottomrule
\end{longtable}
\endgroup
Every entry exceeds six; the ninth has margin \(29/576\).
We justify the entire finite table next.

\paragraph{The \(p\) prefix and its reversed length-three pair.}
\phantomsection\label{part:odd-two:p-three}
Add \(-U-\tau\one\) to the displayed complete primitive.
At \(L=1\), its nonconstant coordinates are
\(\gamma+x(e-E),x(f-C)\). Because \(|e-E|\le x\) and
\(|f-C|\le1\), its full square, including \(-\tau\one\), is at most
\[
 [d^2+2d\theta^2+\theta^4+\theta^3+d]d<2d/3.
\]
The constant cross costs at most
\(\tau^2+2\tau\delta_2=d^2-\delta_2^2\le d^2\).
Its scalar plus mean is at most
\(-4d+2\delta_2+\delta+\delta_3\le-12d/5\).

At \(L=2\), the nonconstant coordinates are
\(1-x^2f-xE,\gamma+x(e-C)\), and the constant is
\(c=\delta-\tau\in[-d,0]\).
The real coordinate is positive and its square is at most \(1-x\);
the imaginary absolute value is at most \(d+x<1\).
Thus its nonconstant square is at most \(d\). Since its mean is
at least \(-\gamma d-\delta_2\), the full square is at most
\(\{1+d(1+2d+2\theta)\}d\le263d/144<15d/8\).
The actual scalar plus mean is at most
\(-5d+3\delta_2+\delta_3\le-18d/5\).

At \(L=3\) the nonconstant coordinates are
\[
 1+\delta-\gamma x-x^2e-xE,\qquad1-x^2f-xC.
\]
They are positive: the real lower bound is
\(1-2\theta/3-\theta^2-\theta=485/3072\).
Their square is at most \((43/36)^2+\theta<15/8\).
The actual scalar plus mean, before the constant \(c=\epsilon-\tau\), is
\[
\begin{aligned}
 &-\{\Gamma-(1-d)\delta\}-2d-g_r
 -\gamma(d-\delta_2)-\int xe\,d\nu\\
 &\quad+\int x^2(e+f)\,d\nu+\int xE\,d\nu\\
 &\le-\tfrac{29}{5}d+2\delta_2+3\delta_3<-179d/40.
\end{aligned}
\]
Transport the same \(c\) by the actual \(p\) curvature;
\(|c|\le29d/16\) and its expense is below \(15d/4\).

For \(L\ge4\) the nonconstant primitive is
\(G_nu+B^nR-U\), and \(c=-hd-\tau\).
The coordinates of \(G_nu-U\) are nonnegative and bounded by
\(1-2x+2x^2,1\); their square is at most one, since
\(-3+8x-8x^2+4x^3<0\) on \([0,\theta]\).
The actual finite tail has norm below \(16\sqrt d/75\), giving
\(\|V\|\le91\sqrt d/75<5\sqrt d/4\).
Its scalar plus mean is at most
\(-[3/5+3h+3-3\theta]d=-159d/40\).
The actual \(p\) curvature pays \(c\), with expense below \(8d/3\).
This proves the first four rows.

\paragraph{The \(p\) prefix and every higher first pair.}
\phantomsection\label{part:odd-two:p-higher}
Add the one actual mask \(V_*\). At length one use
\(\|R+V_*\|\le47\sqrt d/40\) and
\(-\Gamma+\epsilon-16d/5+\E BV_*\le-247d/60\).
At length two the whole primitive before \(V_*\) is
\(u+BR+\delta\one\). Its square divided by \(d\) is at most
\[
 1+\theta(d+\theta)^2+(7/6)d(d+\theta)+(49/144)d<(13/10)^2.
\]
This includes the absolute constant cross, using
\(|\delta|\le7d/12\) and
\(|\E(u+BR)|\le(d+\theta)d\).
The scalar plus mean is at most
\(-[2-\theta^2-1/12+16/5]d=-76057d/15360\).
At length three the nonconstant coordinates are
\(1+\delta-\gamma x-x^2e,1-x^2f\). Their square is below
\(15d/8\), and including \(\epsilon\one\) costs at most
\(195d/256\). Their full square is therefore below
\(675d/256<(13/8)^2d\). Adding \(V_*\) gives norm at most
\(2\sqrt d\). The scalar plus mean is at most
\(-[4/5+2+16/5-1/10-1/12]d=-349d/60\), since
\(-\int xe+\int x^2(e+f)\le-\delta_2+3\delta_3<d/10\).
For the long rows, pay \(c=-hd\) by \(p\) curvature.
The remaining norm is below
\((1+16/75+3/8)\sqrt d<8\sqrt d/5\), its scalar plus mean
is at most \(-2549d/480\), and its curvature expense is below
\(d/3\). This proves the next four rows.

\paragraph{The \(W_2\) prefix and the exchanged length-three pair.}
\phantomsection\label{part:odd-two:w2-three}
At length two the nonconstant coordinates are
\(xE-x^2f,\gamma+xe+xC\), the constant is
\(c=\delta+\beta\), and the scalar is
\(-\Gamma+\epsilon-g_e\).
The scalar plus mean is at most \(-51d/16\), using
\(D_2-\int xE\le-\delta_2+2\delta_3\le0\).
The nonconstant square is at most \(20d/27\), while
\(0<c\le19d/12\) and its mean is nonpositive. Thus the full
square is at most \(20d/27+361d/432=227d/144\).

At length three the coordinates are
\(\delta-\gamma x-x^2e+xE,1-x^2f+xC\), and
\(c=\epsilon+\beta\in[-7d/32,d]\).
Their real absolute value is below \(5/8\); their imaginary
coefficient lies between \(5/6\) and \(17/12\).
The nonconstant square is below \(5d/4\). If \(c\ge0\), its
constant correction is nonpositive. If \(c<0\), that correction
is at most \(2051d^2/3072<d/4\). The full square is therefore
below \(3d/2\). Its exact scalar plus mean is at most
\(-[4/5+1+35/16]d=-319d/80\); the two negative first moments
absorb the two positive second moments before this estimate.

For the long rows the nonconstant primitive is
\((G_n-1)u+B^nR+U\), and \(c=\beta-hd\in[d/16,15d/32]\).
Before the tail, its coordinate bounds are
\([-x^2,x^2]\) and \([1-x,1+x]\), giving square below one.
Its full norm before \(c\one\) is below \(5\sqrt d/4\).
The imaginary coefficient after the actual finite tail is at least
\(1-4\theta/3-\theta^2=901/3072>15/64\).
Consequently \(c(c+2\E V)\le0\), so the full constant does not
increase this square. The scalar plus mean is at most \(-541d/160\).
This proves all three exchanged-pair rows, including \(G_1\).

\paragraph{The \(W_2\) prefix and every higher first pair.}
\phantomsection\label{part:odd-two:w2-higher}
At length two the complete vector is \(BR+\delta\one+V_*\).
Its norm is at most \((3/5+7/20+3/8)\sqrt d=53\sqrt d/40\),
and its scalar plus mean is at most \(-79d/20\).
At length three the nonconstant coordinates before \(V_*\) are
\(\delta-\gamma x-x^2e,1-x^2f\).
The real absolute value is below \(1/2\): bound its two signs
separately by \(7/36\) and \(2\theta/3+\theta^2\).
The nonconstant norm is below \(5\sqrt d/6\).
The actual \(\epsilon\one\) has norm below \(\sqrt d/2\),
so the complete norm is at most \(41\sqrt d/24\).
The scalar plus mean is at most \(-289d/60\).
For the long rows, keep \(c=-hd\) and the actual \(W_2\) curvature.
Transporting it replaces the primitive by
\((G_n-1+hd)u+B^nR+V_*\). Its norm is below
\((4/5+16/75+3/8+17/96)\sqrt d<8\sqrt d/5\).
The extra scalar correction is
\(-4hd^2+(10/3)h^2d^2\le0\), while the uncorrected scalar
plus mean is at most \(-2069d/480\). The last three rows follow.

\paragraph{The distinct unit-length minimum-two direction.}
\phantomsection\label{part:odd-two:unit}
Here the actual terminal is \(W_2\bar f_t\), rather than
\(\bar f_2W_t\). The finite convolution gives
\[
 W_2+\bar f_t\le2-\beta_t+B^2G_{t-2}u.
\]
Thus \(F\le e^{R^--\beta_t}\), with
\(R^-=(-xe-zxf)u\), and
\[
 TF\le e^{-\beta_t}(q+BR^-+\delta^-),\qquad
 \delta^-=\int xf\,d\nu+(3/2)\int x^2(e^2+xf^2)\,d\nu.
\]
The actual \(f_2\) cancels \(u\), leaving primitive \(u+R^-\).
For \(a=3\), add the positive length-three pair with primitive
\((-xE+zxC)u+\tau\one\). The full nonconstant coordinates are
\(1-x(e+E),x(C-f)\). The true constant correction is at most
\(\tau^2+2\tau\delta_2\le d^2\le d/3\).
Keeping all terms in the same original spectral integral bounds the
remaining fourth exponent by \(\int P(x)\,d\nu\), where
\[
 P=-63/4+7x+36x^2-(95/4)x^3+17x^4.
\]
Indeed use \(-4e+4xe\le-4(1-x)^2\), \(f,E\le1\),
\(|C-f|\le1\), and
\(0\le1-x(e+E)\le1-2x+2x^2\) in the exact fourth expansion.
The polynomial increases on \([0,5/12]\), where its endpoint is
\(-20189/2592\). The paid constant gives loss
\(16517/2592>6\).
Its nonconstant square is at most \((1+\theta^3)d<(7/6)^2d\),
and \(\tau\le3\sqrt d/5\); its complete norm is therefore
less than \(2\sqrt d\), as needed for later pairs.
For \(a\ge4\), the base primitive has norm at most \(\sqrt d\),
and its scalar plus mean is at most \(-71d/128\): its integrand
is bounded by the increasing cubic
\(-2+3x+x^2/2+3x^3/2\) at \(5/12\).
The actual root mask then gives complete norm at most
\(11\sqrt d/8\) and loss \(25537/3840>6\).
Both calculations include \((e,f)=(1,0)\) when \(t=3\).

\paragraph{All later pairs, decorations and the outer length.}
\phantomsection\label{part:odd-two:outer}
Every complete or transported vector above has norm at most
\(2\sqrt d\) on the whole host domain; its already paid curvature
constant is unchanged by a later root pair. The entire ordered list
of \(k\ge1\) later pairs is one finite mask, with scalar loss
\(16kd/5\), norm at most \(3\sqrt d/8\), and signed mean at most
\(d/12\). Adding it to the same certificate costs at most
\[
 \left[\tfrac13+\tfrac{17}{4}
       \{2(2)(3/8)+9/64\}\right]d.
\]
The resulting gain is
\(21097/3840+64(k-1)/5>z_k\), proving the first claim.
Equal-pair products lie in \([0,1]\) in their actual compatible
directions, and positivity restores them at both original vertices.

If \(\mathcal M\) is any original positive operator
and \(r=\mathcal M\one\), row Cauchy and H\"older give
\[
 I_q(\mathcal MH)^3\le\E(r^4/q)\|\mathcal M\|^4(\E|H|^4)^2.
\]
This uses the original \(q\), which is positive on every retained atom.
The operator bounds and path rows of
Sections~\ref{proof:Q5-C084} and~\ref{proof:Q5-C164} give exactly
the costs \(204/19,10,0,7,6\) in the five stated length ranges:
at \(M=1\) the two costs are \(4\) and \(128/19\);
at \(M=2\) they are \(6\) and \(4\);
at \(M=3\) both are nonpositive; at \(M\ge4\) the operator
norm is at most one and the path-row costs are \(7,6\).

Since \(12-204/19=24/19\), it proves the second claim for every
independent outer length.
When \(d=0\), leaf integration makes every actual message one,
even if \(B\ne0\).  All finite indices,
continuous endpoints and arm directions have been retained.
\end{proof}
